% Source: round6/attainment/traps.tex
\section{Collision ascent has nonglobal strict local maxima}
\label{s13:sec:attainment-traps}
Let $\mathcal P_d$ be the polytope of probability measures on
$\{-1,1\}^d$ with all coordinate means and distinct pair correlations
zero, and put $C(w)=\sum_xw_x^2$. Finite-field character constructions
of pairwise independent distributions are classical; see Luby--Wigderson,
\emph{Pairwise Independence and Derandomization}, Sections 1.1--1.2,
\url{https://www.math.ias.edu/~avi/PUBLICATIONS/LubyWig2005.pdf}.
The standard atom bound $w_x\le1/(d+1)$ and Hadamard equality case
are discussed by Gavinsky--Pudl\'ak,
\url{https://users.math.cas.cz/~pudlak/pairwise.pdf}.
The results below are self-contained landscape certificates.

\begin{lemma}[Local maximum criterion]\label{s13:lem:attainment-criterion}
A point $w$ of a finite polytope $P$ is a local maximum of the squared
Euclidean norm if and only if $w$ uniquely maximizes the linear
functional $v\mapsto w\cdot v$ over $P$. In this case it is a strict
local maximum and a vertex. In particular, every uniform distribution
that is a vertex of $\mathcal P_d$ is a strict local maximum of $C$.
\end{lemma}
\begin{proof}
If $v\ne w$ has $w\cdot(v-w)\ge0$, then
$\|w+t(v-w)\|^2-\|w\|^2=2t w\cdot(v-w)+t^2\|v-w\|^2>0$
for every $t>0$. Conversely, assume strict linear exposure. Over the
finitely many other vertices $u$, the ratio
$-w\cdot(u-w)/\|u-w\|$ has positive minimum $a$. For any convex
combination $v$ of vertices, linearity and the triangle inequality give
$w\cdot(v-w)\le-a\|v-w\|$. The quadratic increment is therefore
negative when $0<\|v-w\|<2a$.
For a uniform distribution on support $S$ of size $m$,
$w\cdot(v-w)=-\sum_{x\notin S}v_x/m$. A vertex has no other feasible
measure supported on $S$: a second such measure would permit a small
continuation from $w$ in the opposite direction inside its positive
support, contradicting extremality. This proves strict exposure.
\end{proof}

\begin{theorem}[Explicit four- and five-bit traps]
\label{s13:thm:attainment-small-traps}
The following measures are nonglobal strict local maxima of $C$:
\begin{enumerate}
\item On four bits, writing $k$ for the number of negative coordinates,
assign mass $1/6$ at $k=0$, mass $1/12$ to every point with $k=2$ or
$k=3$, and mass zero otherwise.
\item On five bits, assign mass $1/6$ at $k=0$, mass $1/12$ to each
of the ten points with $k=3$, and zero otherwise.
\end{enumerate}
Both have $C=7/72$, whereas feasible measures with $C=1/8$ exist.
\end{theorem}
\begin{proof}
Exchangeability reduces balanced pairwise independence to
$\mathbb E k=d/2$ and $\mathbb E k(k-1)=d(d-1)/4$, verified directly
for both measures. Their collision is $1/36+10/144=7/72$.
For four bits the quadratic function
\[
 q_4(k)=\frac{k^2-5k+12}{72}
\]
equals $w_x$ on its support and exceeds it by $1/9$ elsewhere.
For five bits use
\[
 q_5(k)=\frac1{12}+\frac{(k-3)^2}{108};
\]
it equals $w_x$ on the support and is positive off the support.
The expectation of either quadratic is fixed by the moment constraints,
so $w\cdot v\le w\cdot w$, with equality only if $v$ uses no new atoms.
For four bits the eleven support moment columns are linearly independent;
the accompanying exact checker verifies this using rational elimination.
For five bits this independence is explicit: $1-k/3$ restricts to the
indicator of the all-positive atom. For each two-element set $\{i,j\}$,
the polynomial $\frac14(1+x_i)(1+x_j)-(1-k/3)$ restricts to the
indicator of the weight-three atom with positive coordinates $i,j$.
Thus every supported weight is fixed by the degree-two moments.
The criterion lemma proves strict local maximality. Uniform laws formed
from four or five distinct nonzero linear characters on three independent
fair bits have eight distinct atoms and collision $1/8$; for example
use masks $(1,2,3,4)$ and $(1,2,3,4,5)$ in $\mathbb F_2^3$.
\end{proof}

\begin{theorem}[A trap at admissible Hadamard order eight]
\label{s13:thm:attainment-admissible-trap}
Let $z$ be uniform on $\mathbb F_2^4$, and take seven sign bits
$x_b(z)=(-1)^{b\cdot z}$ with
\[
 B=\{1,2,3,4,5,8,10\},
\]
where masks are written in binary. Their law $w\in\mathcal P_7$ is
a strict local maximum with $C(w)=1/16$, although the global maximum
is $1/8$. In particular, every feasible $v\ne w$ with
$\|v-w\|_\infty<1/136$ satisfies $C(v)<C(w)$.
\end{theorem}
\begin{proof}
Distinct nonzero characters are unbiased and pairwise orthogonal.
The masks include $1,2,4,8$, so the sixteen atoms are distinct.
Moreover
\[
 \{0\}\cup B\cup\{b+b':b,b'\in B,\ b\ne b'\}
 =\mathbb F_2^4.
\]
Therefore the degree-at-most-two monomials restricted to these atoms
span all sixteen Walsh characters. They determine every supported
weight uniquely, proving the uniform measure is a vertex.

For the stated radius, choose for each character one degree-at-most-two
monomial representing it on the support $S$. Walsh inversion gives
polynomials $L_s$ with $L_s(t)=\mathbf1_{s=t}$ on $S$ and
$|L_s(x)|\le1$ on the entire seven-bit cube: each is a sum of sixteen
signed monomials with coefficients $1/16$.
Put $\delta=v-w$ and $\epsilon=\sum_{x\notin S}v_x$.
The fixed moments give
$|\delta_s|=|\sum_{x\notin S}v_xL_s(x)|\le\epsilon$.
Hence $\|\delta\|_1\le17\epsilon$. Since
$w\cdot\delta=-\epsilon/16$, we obtain
\[
 C(v)-C(w)\le-\epsilon/8+
             17\epsilon\|\delta\|_\infty<0.
\]
If $\epsilon=0$, uniqueness forces $v=w$.
Finally, the seven nonzero characters of $\mathbb F_2^3$ give a
uniform eight-atom law with collision $1/8$. No law can exceed that:
for each support atom $s$, the orthonormal functions $1,x_1,\ldots,x_7$
form an isometry into weighted atom space, whose projection diagonal
is $8w_s\le1$. Thus $C\le\max_s w_s\le1/8$.
\end{proof}
These traps concern the full feasible polytope, with arbitrary feasible
perturbations, not only a selected set of coordinate moves. A continuous
nondecreasing collision-ascent path cannot leave a strict local trap.
A finite nonlocal jump can escape. Therefore local collision ascent is
not an all-order attainment proof, even in a dimension where a Hadamard
matrix is already known to exist. No new existence result is asserted.

\paragraph{Exact small enumeration.}
For three bits the feasible Walsh form is
$w_x=(1+a x_1x_2x_3)/8$, $|a|\le1$; its only local maxima are the two
global parity laws. For four bits the checker enumerates all vertices
in the five free Walsh coefficients (four triples and one quadruple).
It finds exactly 26: ten with collision $1/8$ and sixteen with $7/72$.
Every vertex strictly exposes itself: the largest $w\cdot(v-w)$ among
other vertices is negative. The explicit proofs above do not depend
on treating numerical optimization as a certificate.

\begin{theorem}[Arbitrarily poor traps at admissible orders]
\label{s13:thm:attainment-infinite-traps}
For every integer $s\ge2$, put $h=2^s$. There is a uniform vertex
of $\mathcal P_{2h-1}$ supported on $h^2$ atoms. It is a strict
local maximum with collision $1/h^2$, whereas the global maximum
is $1/(2h)$. Their ratio $2/h$ tends to zero. A sufficient strict
local radius in the infinity norm is $2/[h^2(h^2+1)]$.
\end{theorem}
\begin{proof}
Write $G=\mathbb F_2^s\oplus\mathbb F_2^s$. Take all $h-1$
nonzero characters supported on the first summand, all $h-1$
nonzero characters supported on the second summand, and one character
nonzero on both summands. These $2h-1$ characters define sign bits
on a uniform seed $z\in G$. They include a basis of the dual group,
so the $h^2$ atom vectors are distinct. Every group character is
constant, one selected character, or a product of two selected
characters: a mixed character $(a,b)$ is the product of $(a,0)$ and
$(0,b)$. Thus the restricted degree-two moments span all functions
on the support, making the uniform measure a vertex.
The preceding criterion proves strict local maximality. The same
Walsh interpolation proof with $m=h^2$ yields
$\|\delta\|_1\le(m+1)\epsilon$ and
\[
 C(w+\delta)-C(w)
 \le\epsilon\big[-2/m+(m+1)\|\delta\|_\infty\big],
\]
which gives the radius. At this number of bits, all nonzero characters
on $\mathbb F_2^{s+1}$ give $2h$ equal-weight atoms and collision
$1/(2h)$. The atom bound proves global optimality.
\end{proof}
The construction uses two independent seed blocks but includes their
complete sets of characters and a mixed character as observed bits;
it is not an independent product on two disjoint observed bit blocks.
The latter distinction is essential: a product whose support exceeds
the number of moment equations cannot be a local maximum, since it
has a nonzero two-sided supported feasible direction.

\begin{corollary}[A certified collision barrier]
\label{s13:cor:attainment-barrier}
For each uniform character trap above with support size $m$, every
continuous feasible path from the trap to a global maximizer passes
through a measure of collision at most
\[
 \frac1m-\frac1{m^2(m+1)}.
\]
\end{corollary}
\begin{proof}
The interpolation bound implies $\|\delta\|_\infty\le\epsilon$:
it holds on the support, and each off-support mass is at most their
total $\epsilon$. At infinity distance $r=1/[m(m+1)]$, the preceding
inequality gives $C(w+\delta)-1/m\le-\epsilon/m\le-r/m$.
Every global maximizer is uniform on $n=d+1<m$ atoms: the atom bound
$w_x\le1/n$ and $C=1/n$ force each positive weight to equal $1/n$.
It therefore omits at least one trap atom, putting it at infinity
distance at least $1/m>r$. A continuous path must cross the radius-$r$
boundary.
\end{proof}


\paragraph{Formalization of the local maximum criterion.} The theorem \texttt{finite\_\allowbreak polytope\_\allowbreak local\_\allowbreak max\_\allowbreak criterion} proves the stated equivalence for any finite convex hull in a real inner-product space, together with strict locality and the extreme-point conclusion. The theorem \texttt{uniform\_\allowbreak cube\_\allowbreak vertex\_\allowbreak strict\_\allowbreak local\_\allowbreak collision} proves the uniform-vertex consequence for the full Boolean-cube probability, mean, and pair-correlation polytope. Its independent proof derives a strict linear margin on the closed feasible tangent cone by compactness of its unit section. The moment-matrix representation of the cube constraints is proved explicitly; no polytope decomposition or local-maximality oracle is assumed.

\paragraph{Formalization of the infinite character family.} The theorem \texttt{arbitrarily\_\allowbreak poor\_\allowbreak character\_\allowbreak traps} proves all finite-parameter conclusions, including the actual full-polytope vertex, strict local maximum with the displayed radius, and the attained global comparison. The theorem \texttt{character\_\allowbreak trap\_\allowbreak relative\_\allowbreak collision\_\allowbreak limit} proves that the ratio tends to zero. The observed-coordinate index type is proved to have exactly $2h-1$ elements. The proof constructs bounded Walsh interpolation functions, rather than assuming their existence. The separate specified order-eight mask example and the barrier statement covering every preceding example have their own coverage records.

\paragraph{Formalization of the order-eight example and barriers.} The theorem \texttt{order\_\allowbreak eight\_\allowbreak character\_\allowbreak trap} uses exactly the binary masks $1,2,3,4,5,8,10$, proves the sixteen-atom support and all moment constraints, and establishes the explicit radius $1/136$ against every feasible perturbation. Theorems \texttt{character\_\allowbreak trap\_\allowbreak global\_\allowbreak path\_\allowbreak barrier} and \texttt{order\_\allowbreak eight\_\allowbreak global\_\allowbreak path\_\allowbreak barrier} cover the infinite family and this specified example, respectively. The interpolation estimate is proved algebraically, and an intermediate-value argument produces the actual crossing point on each continuous path.

\paragraph{Formalization of the two nonuniform small traps.} The theorems \texttt{four\_\allowbreak bit\_\allowbreak nonglobal\_\allowbreak trap} and \texttt{five\_\allowbreak bit\_\allowbreak nonglobal\_\allowbreak trap} prove the stated weight formulas, exact moments, collision $7/72$, strict locality in the full feasible polytope, and larger-collision witnesses. Their rational dual and support-interpolation certificates are generated by an original Python standard-library script and checked by Lean kernel reduction. The proof then transfers the certificates to real laws and applies the proved tangent-cone margin theorem; generation is not trusted as a proof oracle.
